% !PRINTED 665
\weberpage{665}

%% [running head, not body text -- commented out 2026-08-07] Page furniture (running
%% head, folio, signature) is excluded from this edition by standing policy.
%% \S.\,182. \hfill \gesperrt{Beweis der H\"ulfss\"atze.} \hfill 665

Um $S_n$ zu bilden, multipliciren wir diese Gleichung mit $t_1'$ und
summiren. Dadurch ergiebt sich

\[
S_n \equiv \sum^{t_1'} t_1' \sum^{t_2} E\left(\frac{n t_2}{m}\right)
+ \sum^{t_1} t_1' E\left(\frac{n t_1}{q}\right) \pmod q.
% [source sic: the denominator m repeats the source's m-for-m' error from the preceding leaf]
\]

Darin ist nun $\sum t_1' = \sum t_1 \equiv 0 \pmod q$, weil die $t_1$
paarweise die Summe $q$ ergeben, und es folgt:

\weq{15}{S_n \equiv \sum^{t_1} t_1' E\left(\frac{n t_1}{q}\right) \pmod q.}

Die Summe, die hier auf der rechten Seite steht, ist ebenso gebildet,
wie $S_n$ selbst in dem Falle, wo $m$ eine Primzahl ist. F\"ur diesen
Fall ergiebt sich aber

\[
\frac{n t_1}{q} = E\left(\frac{n t_1}{q}\right) + \left(\frac{n t_1}{q}\right)
\]
\[
\frac{(q-n) t_1}{q} = E\left(\frac{(q-n) t_1}{q}\right) + \left(\frac{(q-n) t_1}{q}\right),
\]

und daraus durch Addition

\[
t_1 = E\left(\frac{n t_1}{q}\right) + E\left(\frac{(q-n) t_1}{q}\right)
+ \left(\frac{n t_1}{q}\right) + \left(\frac{(q-n) t_1}{q}\right),
\]

und dabei kann die Summe der beiden positiven echten Br\"uche, die
doch eine ganze Zahl sein muss, nur den Werth 1 haben. Es folgt also

\[
E\left(\frac{n t_1}{q}\right) + E\left(\frac{(q-n) t_1}{q}\right) = t_1 - 1,
\]

und daraus durch Multiplication mit $t_1'$ und Summation \"uber alle
Werthe von $t_1$ von 1 bis $q - 1$

\[
\sum^{t_1} t_1' E\left(\frac{n t_1}{q}\right)
+ \sum^{t_1} t_1' E\left(\frac{(q-n) t_1}{q}\right) \equiv -1 \pmod q.
\]

Folglich k\"onnen sicher nicht beide Summen auf der linken Seite durch
$q$ theilbar sein. Dann zeigt der Ausdruck (15), dass von den beiden
Summen $S_n$ und $S_{q-n}$ gewiss eine durch $q$ untheilbar ist.

Dies aber war zu beweisen.
